\chapter{Lua API}

Lua 5.4 scripts are loaded when Lua is available. Scripts return a table and may implement lifecycle functions.

\section{Lifecycle}

Supported lifecycle functions are:

\begin{lstlisting}[language={[5.2]Lua}]
function Script:on_create() end
function Script:on_start() end
function Script:on_update(dt) end
function Script:on_fixed_update(dt) end
function Script:on_destroy() end
\end{lstlisting}

HUD button scripts may additionally implement:

\begin{lstlisting}[language={[5.2]Lua}]
function Script:on_ui_hover(event) end
function Script:on_ui_click(event) end
\end{lstlisting}

The UI event contains \code{id}, \code{label}, \code{mouse\_x}, and \code{mouse\_y}. Mouse coordinates are in the HUD file canvas coordinate system.

\section{Services}

The current Lua services include:

\begin{longtable}{>{\ttfamily}p{0.30\textwidth}p{0.58\textwidth}}
Debug & Logging, debug lines, and clearing debug lines. \\
Input & Keyboard, mouse buttons, mouse position, mouse world position, viewport size. \\
Entity & Entity lookup, creation, destruction, and transform position access. \\
Rigidbody2D & Get and set velocity and apply impulses. \\
Physics2D & Overlap box queries. \\
Hud & Imperative HUD text, rects, visibility, colors, groups, and text updates. \\
Save & Simple number and string persistence in save slots. \\
Audio & Play, stop, master volume get and set. \\
Runtime & Quit, physics enable, window mode get and set. \\
Scene & Runtime scene switching by scene id. \\
\end{longtable}

\section{Save API}

The save API is designed for small game settings and simple persistent values.

\begin{lstlisting}[language={[5.2]Lua}]
local volume = Save.get_number("settings", "master_volume", 1.0)
Save.set_number("settings", "master_volume", volume)

local mode = Save.get_string("settings", "window_mode", "windowed")
Save.set_string("settings", "window_mode", mode)
\end{lstlisting}

\section{HUD API}

HUD elements can be generated or updated from Lua:

\begin{lstlisting}[language={[5.2]Lua}]
Hud.text("points", "POINTS: 0", 24.0, 24.0, 3.0)
Hud.rect("panel", 20.0, 20.0, 320.0, 120.0, 0.1, 0.1, 0.2, 0.8)
Hud.set_text("points", "POINTS: 10")
Hud.set_rect("bar", 100.0, 100.0, 200.0, 16.0)
Hud.set_color("menu_tab_sound", 0.24, 0.27, 0.46, 0.94)
Hud.set_visible("points", true)
Hud.set_group_visible("menu_main", false)
\end{lstlisting}

\section{Scene Switching}

\code{Scene.load} requests a runtime scene switch by scene id. The switch is deferred: it is applied between frames by the runtime loop, so it is safe to call from any Lua callback (menu buttons, update, fixed update).

\begin{lstlisting}[language={[5.2]Lua}]
Scene.load("scene://minimal_2d/spiral")
\end{lstlisting}

When a scene switch is applied, the runtime:

\begin{enumerate}
  \item validates and parses the target scene (and its optional HUD) by id from the project scene list;
  \item calls \code{on_destroy} on every active script and releases their script tables;
  \item replaces the active world in place (so existing world references stay valid);
  \item re-loads the project entry script, entity scripts, and HUD button scripts;
  \item calls \code{on_create} and \code{on_start} on the freshly loaded scripts.
\end{enumerate}

Required modules stay cached across a switch, so a shared state module can carry a pending flag (for example an \code{auto\_start} flag) from the button that requested the switch into the new scene's \code{on_start}. Runtime-created entities from the previous scene are discarded with the old world.
